\chapter{元素的身份证}

\begin{kaichang}
你手里多半有一支铅笔。铅笔杆上常印着“HB”或“2B”，中文叫它“铅笔”，英文里也叫 lead pencil——“铅”笔。可是我要告诉你一件有点扫兴的事：铅笔芯里\textbf{一丁点儿铅都没有}。它是石墨和黏土压出来的，而石墨是纯粹的碳。

名字叫铅，实际是碳。那到底由什么说了算：一种元素“是”什么，究竟由什么决定？再看一个更古老的执念：几百年间，无数炼金术士想把灰扑扑的铅变成金灿灿的金，烧炉子的、加药的、念咒的，前赴后继，全部失败。他们是手艺不行，还是这件事在根本上就不可能？

先写下你猜的答案。这一章读完，我们会给每一种元素发一张“身份证”，然后一眼看穿这两个问题。
\end{kaichang}

\section{化学反应翻不动的底牌}

先看现象。铁钉会生锈，铜锅会发绿，白糖烧久了变黑——物质的样子千变万化。十七、十八世纪的化学家把各种物质烧、煮、溶解、通电，折腾了两百多年，发现一件怪事：有些东西无论怎么折腾，都变不成别的东西。

把铁烧红、打成片、磨成粉、溶进酸里，铁永远是铁；让它和氧气、硫、氯分别反应，生成物里照样找得到铁。把金拉成比头发细得多的丝，用王水溶掉，再想办法析出来——它还是金。拉瓦锡那一代人正是靠这份“怎么折腾都不变”的清单，反过来定义了元素：元素就是化学反应无法分解、也无法互相转化的那些基本成分。

换句话说，化学反应像洗一副扑克牌：牌可以换组合、换搭档，但牌本身一张张都还在。烧一块木炭，碳并没有“消失”，它只是换了个搭档，跑进了空气里（第 26 章会用账本把这件事算得清清楚楚）。

生活里到处都是这份清单的影子。金戒指戴了十年还是金的，顶多磨掉一层皮；加碘盐里的碘，溶进汤里、吃进肚里，碘还是碘；你吸进一口空气，氮气在肺里转了一圈又原样呼出来——身体根本“拆不动”它。古人不懂这套清单，才会把“点石成金”当成只是技术问题。那么问题就被推到下一层：牌面底下，到底是什么东西锁死了一种元素的身份？

\section{质子数：身份证号码}

第 6 章已经见过原子内部的三位主角：带正电的质子和不带电的中子，挤在比原子小几万倍的原子核里；核外是带负电的电子。这三种粒子，哪一种决定元素的身份？

一个一个排除。电子最可疑，也最先出局：钠光灯里的钠原子和食盐里的钠离子，电子数明明不一样（一个 11 个，一个 10 个），化学家却说它们都是“钠”。看来电子数说了不算。中子呢？马上你会看到，中子数不同的原子也可以是同一种元素。剩下质子——它才是身份证号码：

\textbf{一种元素，就是所有原子核里质子数相同的一类原子的总称。}质子数有一个专门的名字，叫\textbf{原子序数}。氢的原子序数是 1，碳是 6，氧是 8，铁是 26，金是 79。全世界每一个带着 6 个质子的原子核，不管外面围着几个电子、核里有几个中子、此刻是藏在你体内还是镶在钻石上，它都是碳。而铅有 82 个质子，金有 79 个——中间隔着 3 个质子，这就是铅变不成金的全部原因。

这个结论今天看来理所当然，可“原子序数”最初只是元素在表格里的排位号，谁也不敢说它对应原子里的什么实物。1913 年，英国青年物理学家莫塞莱用 X 射线轰击各种元素，发现每种元素发出的 X 射线频率都严格随一个整数台阶移动——这个整数正是原子序数，而且可以直接解释为原子核所带的正电荷数。原子序数从此有了物理含义：它就是核里的质子数。顺带一提，这件事还了门捷列夫一个公道：他当年坚持按“性格”把碲排在碘前面，与原子量顺序相反，被不少人嘲笑；原子序数一量，他的排法分毫不差（第 10 章还会讲到这段故事）。莫塞莱 1915 年战死在加里波利前线，年仅 27 岁，科学史上常有人感叹：他是最令人痛惜的“没有拿到诺贝尔奖”的人之一。

为什么偏偏是质子数说了算？道理藏在因果链上：核里的质子带正电，正电的多少决定了一个中性原子能“抓住”多少个电子；而原子的化学行为——爱跟谁反应、怎么反应——几乎全部由核外电子的数目和排布决定（第 9 章和第 11 章会一层层拆开这件事）。所以质子数在最上游：它定了电子的家底，电子的家底定了性格。中子躲在核里不带电，对电子毫无发言权；电子自己则进出自如。一场戏看下来，只有质子从没换过演员。

\section{多几个中子：同位素}

身份证号码锁死的是质子数。可同一个号码下，原子核还可以有不同的“体重”——因为中子数可以变。

拿碳来说。质子数与中子数之和有个名字，叫\textbf{质量数}。绝大多数碳原子的核里是 6 个质子加 6 个中子，质量数 12，写成碳-12；但每大约 100 个碳原子里，就有 1 个多带一个中子，是碳-13；还有极其稀少的碳-14，核里是 6 个质子加 8 个中子。质子数都是 6，所以它们全都是碳，化学行为几乎一模一样，只是核的体重不同。这种\textbf{质子数相同、中子数不同的同类原子}，互称\textbf{同位素}——希腊语里“同一个位置”的意思，因为它们在周期表的同一格里（周期表的全貌，第 10 章会展开）。

氢是同位素差别最夸张的元素。普通氢的核里只有 1 个质子；氘多 1 个中子，体重直接翻倍；氚多 2 个中子，体重变成三倍。翻倍是什么概念？碳-13 比碳-12 只重约 8\%，而氘比氢重整整 100\%。所以含氘的水——“重水”——密度比普通水大约 10\%，冰点也高了近 4 度：普通水在 $0\,^{\circ}\mathrm{C}$ 结冰，重水在约 $3.8\,^{\circ}\mathrm{C}$ 就结冰。体重不同还会拖慢化学反应的速度，这叫同位素效应；生命体内的酶对轻重不同的同位素有一点点“偏心”，科学家正是靠测量这种偏心，推断古人主要吃什么、古气候是冷是暖。

绝大多数同位素是安稳的：碳-12、碳-13、氧-16、氘，几十亿年都不会自己变一下。少数同位素的核里质子中子配比不佳，会在某个时刻自发改变——碳-14 和氚都在这份名单上。安稳的同位素早已渗入日常：重水在某些核反应堆里是正式的工作人员（干什么用，第 8 章揭晓）；冰川深处的冰芯里，氧-18 与氧-16 的比例随远古气温起伏，是气候学家最信赖的“温度计”之一。至于不稳定的那些同位素怎么变、变完成了什么，先卖个关子，那是第 8 章的主角。

\section{多几个电子：离子还是它自己}

再说电子。原子整体不带电时，电子数恰好等于质子数。但电子待在原子最外层，是三种粒子里最容易“进出门”的。钠遇到氯，钠原子的一个电子转移到氯那边，钠变成带正电的 \chem{Na^+}，氯变成带负电的 \chem{Cl^-}。带电的原子或原子团，叫\textbf{离子}。

请盯紧这一步里什么变了、什么没变：电子数变了，整体电荷变了，化学性格翻天覆地——银白色、遇水就剧烈反应的金属钠不见了，黄绿色、有毒的氯气也不见了。可是，两边的原子核一个粒子都没动：11 个质子还是 11 个质子，17 个还是 17 个。所以 \chem{Na^+} 仍然是钠，\chem{Cl^-} 仍然是氯——它们只是同一种元素换了“带电状态”。这就是为什么氯化钠（食盐）可以安全入口：盐里的氯和毒气里的氯是同一种元素，但“同一种元素”从来不等于“同一种物质”（第 10 章在食盐与消毒水的例子里也强调过这一点）。

你在超市货架上随时能验证这个道理。运动饮料的配料表写着“钠、钾”，没有任何厂家敢往里面放金属钠和金属钾——那会出人命；瓶子里装的其实是钠离子和钾离子，是这两种元素“交出电子之后”的安稳形态。补钙片里的“钙”也不是灰白色的金属钙块，而是碳酸钙里的钙离子。商家嘴上说的是元素，瓶里装的是离子，消费者稀里糊涂，化学家分得清清楚楚。

一张小表收个尾：

\begin{table}[htbp]
\centering
\begin{tabular}{lll}
\toprule
改变的东西 & 元素身份是否改变 & 结果 \\
\midrule
电子数 & 不变 & 得到离子，化学性质大变 \\
中子数 & 不变 & 得到同位素，化学性质几乎不变 \\
质子数 & 改变 & 变成另一种元素，化学反应做不到 \\
\bottomrule
\end{tabular}
\end{table}

\section{把身份证写成符号}

到这儿，该让符号登场了。元素用一两个字母表示：\chem{C} 是碳，\chem{O} 是氧；有些符号取自拉丁名，比如铁写作 \chem{Fe}，来自 ferrum，所以和中文名对不上。完整版的身份证把号码和体重都标上：左下角写原子序数（质子数），左上角写质量数，例如

\[{}^{12}_{6}\chem{C} \qquad {}^{14}_{6}\chem{C} \qquad {}^{2}_{1}\chem{H}\]

读法很直白：${}^{14}_{6}\chem{C}$ 就是“6 号元素、质量数 14”，也就是碳-14。中子数不用写——用左上角减去左下角就行：$14-6=8$ 个中子。离子则把净电荷标在右上角：\chem{Na^+}、\chem{Cl^-}、\chem{Ca^{2+}}，数字是失去或得到的电子数。

\section{给原子称体重}

\begin{zhengju}
刚才说“绝大多数碳是碳-12，约 1\% 是碳-13”——这种话科学家是怎么知道的？总不能把原子一个个数出来看。答案是一台专门给原子“称体重”的仪器，而它的来历里有一段真实的悬案。

十九世纪末，化学家已经测出各种元素的相对原子质量，结果发现一个别扭的现象：很多元素的原子质量不是整数。比如氖，怎么测都在 20.2 左右——如果每种元素真像道尔顿设想的那样，由完全相同的一类原子构成，这个 0.2 是哪来的？有人怀疑是测量误差，有人怀疑天然元素本身就“不纯”，还有人干脆怀疑原子论。

1913 年，汤姆孙让一束氖离子穿过电场和磁场。按旧想法，所有氖离子质量相同，应该在底片上只打出一条抛物线——结果打出了两条：一条对应质量 20，另一条更暗的对应质量 22。同一种元素，两种体重。他的学生阿斯顿断定这仪器大有可为，1919 年造出了第一台真正精确的\textbf{质谱仪}。它的思路分三步：先把原子变成带电的离子；再用电场把离子加速；最后让它们穿过磁场——磁场会把运动的带电粒子掰弯，而轻的粒子比重的更容易被掰弯，于是不同质量的离子落在不同位置，就像轻重不同的球滚过一阵侧风。阿斯顿用这台仪器查了一大批元素，确认氖是约 90\% 的氖-20 和约 10\% 的氖-22 的混合物，加权平均恰好是 20.2——非整数原子质量的悬案就此告破：错的不是原子论，而是“每种元素只有一类原子”这个隐藏假设。阿斯顿因此获得 1922 年诺贝尔化学奖，而他最重要的结论朴素得可爱：天平上称出来的，从来都是混合物的平均体重。
\end{zhengju}

今天，质谱仪早已走出实验室，成了化学界最忙碌的“秤”。医院里，新生儿足跟的一滴血，用质谱筛查几十种先天代谢病；体育赛场上，它从运动员的尿样里分辨出十亿分之一量级的违禁药物；机场安检口擦一擦你的行李，试纸送进去的就是一台小型质谱仪，专门闻爆炸物的“气味”；农产品检测站里，它检查蔬菜上的农药残留。还有一个和本章直接相关的医学应用：碳-13 呼气试验。医生让你服下一粒含碳-13 标记尿素的胶囊——碳-13 是完全稳定的同位素，毫无放射性——如果胃里有幽门螺杆菌，这种细菌会把尿素拆开，你呼出的气体里就会出现带碳-13 的二氧化碳，仪器一称便知。拿“体重不同的同一种元素”当追踪标签，这是同位素最优雅的用法之一。

\begin{shiyan}
在家可以做一个质谱仪的“桌面类比”。材料：一台吹风机、一个乒乓球、一颗玻璃弹珠（或实心金属球）、一张平整的桌面、一支粉笔。步骤：①在桌面一端用粉笔画一条直线，让球沿这条线自然滚动；②把吹风机固定在桌子侧面，垂直于直线吹出稳定的风，做成一个“风区”；③先后让乒乓球和玻璃弹珠以差不多的速度滚过风区，观察各自的轨迹偏离直线多少。观察点：哪个球被“掰弯”得更厉害？你会发现轻的乒乓球明显偏转，重的弹珠几乎走直线——质谱仪里，轻的离子偏转大、重的离子偏转小，就是这个道理。两点提醒：这只是类比，真实的离子是被电场加速、被磁场偏转，没有谁在吹它们；安全上，吹风机不要对着人或易燃物吹，不要用湿手操作电器。
\end{shiyan}

\section{平均体重：相对原子质量}

有了质谱仪，元素表格里那些带小数的原子质量就不神秘了。化学家规定：取一个碳-12 原子的质量为标准，把它的 $\frac{1}{12}$ 定为 1 份；某种原子的质量相当于多少份，就是它的相对质量。而你在表格里查到的那个数——比如氯的 35.45——是这种元素\textbf{所有天然同位素按含量加权的平均值}，叫\textbf{相对原子质量}。

\begin{qiaoliang}
拿氯来算一笔账。质谱仪告诉我们：天然氯里约 75\% 是氯-35，约 25\% 是氯-37。随机抓一大把氯原子，平均每 4 个里约有 3 个是 35、1 个是 37，所以平均质量约为

\[\frac{35\times 3 + 37\times 1}{4} = \frac{105+37}{4} = 35.5\]

和表格上的 35.45 几乎吻合（剩下那一点差距，是因为同位素的精确质量比整数质量数略小一点点，这里先不计较）。注意一个常被问到的问题：自然界里并不存在质量为 35.45 的氯原子，一个都没有——35.45 是混合物的平均值，就像“平均每个家庭有 1.8 个孩子”，不代表谁家真有 0.8 个孩子。
\end{qiaoliang}

再来一个真实的数量级估算，感受碳-14 有多稀少。一个体重 70 千克的成年人，身体里大约 18\% 是碳，约 12.6 千克。按每 \ud{12}{g} 碳约含 $6\times 10^{23}$ 个原子算（第 5 章的摩尔），你体内的碳原子总数约 $6\times 10^{26}$ 个。活的生物体里，碳-14 占全部碳的比例大约是 $1.2\times 10^{-12}$——每万亿个碳原子里才摊上一个多一点。乘起来，你此刻体内大约有 $7\times 10^{14}$ 个碳-14 原子。七百万亿，听着吓人，可它们稀释在六十万亿亿亿个普通碳原子中间，平均每秒钟只有几千个会自发发生变化。你的身体就这样带着一整个“同位素动物园”在生活，而你毫无感觉。

\section{这张身份证的边界}

每个模型都有边界，这张“身份证”也不例外。

第一，质子数锁定身份，这条规则有明确的适用范围：化学反应。化学反应只动核外电子，连原子核的门都不敲，所以在烧、煮、溶、通电的世界里，铅永远是铅。可一旦能量高到能撬动原子核本身——用粒子轰击、让核聚变或裂变——质子数照样会变，铅甚至真的可以变成金（科学家在粒子加速器里真做出来过，代价是比金子本身贵无数倍的电费和设备）。炼金术士失败，不是身份规则错了，是他们的炉子根本够不着原子核。这条边界里面，藏着第 8 章的全部故事。

第二，身份证只证明“是谁”，不预告“会干什么”。同样是纯碳，排成片层是滑软的石墨，搭成立体网格是坚硬的金刚石，卷成空心球是富勒烯——元素身份一模一样，物质天差地别。同样是 8 个质子的氧，两个原子牵手是供你呼吸的 \chem{O_2}，三个凑一起是刺鼻的臭氧 \chem{O_3}。想知道一种物质的性格，光查身份证不够，还得看原子们怎样连接、排列、交换电子——那正是下一篇“化学键”要拆解的内容。

第三，同位素“化学性质几乎相同”里的“几乎”二字不能丢。对大多数元素，差别小到可以忽略；但对最轻的氢，氘的体重翻倍，差别大到能写进课本——重水能参与普通水的大部分反应，速度却明显更慢，生物若长期只喝重水会出大问题。质量越轻的元素，同位素之间的差别越不能当没看见。

\begin{wuqu}
“同位素就是有放射性的危险原子。”——大概是新闻里“放射性同位素”这个词出现太多次，很多人把两者画了等号。事实是：同位素只是“质子数相同、中子数不同”的同类原子，和放射性没有任何必然关系。你喝的水里有氧-16、氧-17、氧-18 三种同位素，你身体里的碳、钾全都自带天然同位素配比，其中绝大多数稳定得几十亿年都不变一下。会自发衰变的放射性同位素是少数派，而且危险与否还要看剂量和接触方式（第 8 章细讲）。碳-13 呼气试验连孕妇和儿童都能做，正因为它用的是稳定同位素。看到“同位素”三个字先别紧张，先问两句：它稳定吗？剂量多大？
\end{wuqu}

\section{回到铅笔尖}

现在兑现开场的两个承诺。

铅笔为什么叫“铅”笔？十六世纪，英国人发现一处石墨矿，看它灰黑发亮、能在纸上留痕，以为是一种铅矿，“铅笔”的名字就这么以讹传讹地留了下来。可只要查一下身份证：石墨的原子核里 6 个质子，是碳；铅的核里 82 个质子，是铅。名字会骗人，质子数不会。今天反过来的“骗局”也存在：镀金的首饰、仿钻的锆石，肉眼难分，一查原子序数就现形。

铅为什么变不成金？因为两者之间隔着 3 个质子，而一切炉火、药剂、通电都只在电子层面折腾，碰不到原子核。炼金术士几百年失败的真正价值，是逼人类想明白：化学反应改变的是组合，不是身份；想改写身份，得去敲原子核的门。

顺便把你自己的身体也检查一遍：你血液里的铁离子和铁钉是同一种元素，你细胞里的钾离子和遇水就炸的金属钾是同一种元素，你骨骼里的钙离子和石灰里的钙是同一种元素——你是一堆“领过身份证、又各自换了带电状态”的元素，按照精密的连接方式组装起来的。这本书后面八十来章，就是要把这份组装说明书一页页读下去。

\section{门口还站着一位}

这一章我们好几次把碳-14 的故事讲到一半就停住。它核里多出的两个中子，让它无法永远安稳：平均每过五千七百多年，一批碳-14 原子里就有一半完成了变身。这个“一半”不是约定，而是概率——具体到哪一个原子、哪一秒变，谁也预测不了；但几万亿个原子的总体行为，却精确得像钟表。正是这套“个体不可预测、群体极其守时”的规律，让考古学家称得出一块兽骨的年代，让医生能用放射线对付肿瘤，也让人类释放出原子核里沉睡的巨大能量。原子核这扇我们一直假装不存在的门，下一章就去敲开它。

\begin{tiaozhan}
这一段可以跳过，不影响主线。细心的读者可能发现一个逻辑圈套：我们说“相对原子质量以碳-12 质量的 $\frac{1}{12}$ 为标准”，又说“用质谱仪测同位素质量”——那质谱仪的刻度又是拿什么校准的？科学测量里，单位总得有个锚点。2019 年之前，国际单位制里的“千克”由巴黎郊外一个铂铱合金圆柱体定义，“摩尔”则绑在 \ud{12}{g} 碳-12 所含的原子数上。2019 年国际计量大会做了一件大事：把锚点从“某个实物”换成“几个自然常数”，直接规定阿伏伽德罗常数 $N_A = \ud{6.02214076\times 10^{23}}{mol^{-1}}$，千克则由普朗克常数反推出来。从此，碳-12 的摩尔质量不再“恰好”是 \ud{12}{g/mol}，而是一个带有微小不确定度的测量值。定义换了，你称出来的体重一两都不会变——但科学的锚，从一块会磨损的金属疙瘩，换成了宇宙本身的常数。留一道小题：硼的相对原子质量是 10.81，天然硼只有硼-10 和硼-11 两种同位素，能反推出它们各占百分之几吗？（设硼-11 占 $x$，列方程 $11x+10(1-x)=10.81$，解得 $x=0.81$，即约 81\%。）
\end{tiaozhan}

\section*{练一练}

\begin{sikaoti}
\item 画三个方框代表三个原子核，分别画出：6 个质子 6 个中子、6 个质子 8 个中子、8 个质子 8 个中子（质子画成带“$+$”的圆点，中子画成空心圆点）。给每个方框标注它是什么元素、什么同位素（用 ${}^{A}_{Z}\chem{X}$ 记法），再写一句话：哪两个方框里的原子化学行为几乎一样？为什么？
\item 一个氧原子（8 个质子、8 个中子、8 个电子）得到两个电子后变成 \chem{O^{2-}}。请列出变化前后质子、中子、电子各有多少个，并回答：为什么它仍然算“氧”？
\item 天然铜只有铜-63 和铜-65 两种同位素，铜的相对原子质量是 63.55。不查资料，先判断哪种同位素含量更高，并说明理由；再设铜-63 占 $x$，列方程算出具体百分比。
\item 用“吹风机吹乒乓球和弹珠”的类比，向一位没读过本章的朋友解释质谱仪：离子为什么要先带电？为什么轻的偏得多？这个类比在哪个环节会“失灵”？（提示：磁场对静止不动的电荷不起作用，所以离子得先被加速。）
\item 某广告宣称其产品能“把空气中的氧变成负氧离子，更有益健康”。用本章的身份证逻辑分析：\chem{O_2} 得到电子之后，元素身份变了吗？得到的物质还是原来那个物质吗？这条广告混淆了哪两个概念？
\end{sikaoti}
